\documentclass[10pt,a4paper]{article}
\usepackage[T1]{fontenc}
\usepackage{lmodern}
\usepackage{geometry}
\geometry{a4paper,total={170mm,257mm},left=20mm,top=20mm}
\usepackage{amsmath,amssymb,amsthm,mathtools}
\usepackage{microtype,booktabs,array}
\usepackage{xcolor}
\usepackage[numbers,sort&compress]{natbib}
\usepackage{xurl}
\usepackage[colorlinks=true,linkcolor=blue!55!black,citecolor=blue!55!black,urlcolor=blue!55!black]{hyperref}
\hypersetup{pdftitle={Asymptotic Tensor-Rank Sublevels Need Not Be Irreducible},pdfauthor={Wenqi Zhu; Qiyuan Chen; Ping Nie}}
\newtheorem{theorem}{Theorem}[section]
\newtheorem{proposition}[theorem]{Proposition}
\newtheorem{lemma}[theorem]{Lemma}
\newtheorem{corollary}[theorem]{Corollary}
\theoremstyle{definition}
\newtheorem{definition}[theorem]{Definition}
\theoremstyle{remark}
\newtheorem{remark}[theorem]{Remark}
\numberwithin{equation}{section}
\newcommand{\C}{\mathbb C}
\newcommand{\F}{\mathbb F}
\newcommand{\AR}{\widetilde{\mathrm R}}
\newcommand{\AQ}{\widetilde{\mathrm Q}}
\newcommand{\R}{\mathrm R}
\newcommand{\Q}{\mathrm Q}
\DeclareMathOperator{\rank}{rank}
\DeclareMathOperator{\End}{End}
\DeclareMathOperator{\Cent}{Cent}
\DeclareMathOperator{\im}{im}
\DeclareMathOperator{\supp}{supp}
\title{Asymptotic Tensor-Rank Sublevels Need Not Be Irreducible\\[2pt]\large A counterexample in order four and the growth of components}
\author{%
Wenqi Zhu\thanks{Corresponding author.}\\
{\small Mathematical Institute, Woodstock Road, University of Oxford,}\\
{\small Oxford, UK, OX2 6GG}\\
{\small\texttt{wenqi.zhu@maths.ox.ac.uk}}\\[1ex]
Qiyuan Chen\\
{\small Academy of Mathematics and Systems Science,}\\
{\small Chinese Academy of Sciences, Beijing, China}\\[1ex]
Ping Nie\\
{\small David R. Cheriton School of Computer Science, University of Waterloo}\\
{\small\texttt{ping.nie@uwaterloo.ca}}}
\date{}

\renewcommand{\baselinestretch}{1.15}
\setlength{\emergencystretch}{2em}
\begin{document}
\maketitle
\begin{abstract}
The set of complex tensors of asymptotic rank at most $\sqrt[3]{62}$
in format $(3,3,3,2)$ is reducible in the Zariski topology.
This gives a negative answer to the irreducibility question when
quantified over all tensor orders at least three; the order-three
question remains open. A finite Fourier degeneration and Strassen's
asymptotic-spectrum and tight-support theorems give the bound
$\AR(P)^3\leq 74-3\sqrt[3]{81}<62$ for the three-mode permutation
tensor $P$. The scalar centroid of $P$ supplies an ambient open
neighborhood in which the balanced flattening constraints force the
fourth factor to have rank one. A second rank-two tensor then exhibits
two proper relatively closed subsets covering the entire sublevel.
More generally, for every $n\geq3$, every order $k\geq4$ and all
$d_1,\ldots,d_k\geq2$, we exhibit
thresholds $n<B_n<n+1$ such that, for $B_n\leq a<n+1$, the complex
asymptotic-rank sublevel in format $(d_1,\ldots,d_k)$ has at least
$\binom{t_n}{3}+1$ irreducible components, where $t_n=\#\{i:d_i\geq n\}$.
The finite identities and witness geometry admit exact arithmetic
checks and a Lean certificate; the asymptotic and Zariski arguments
are given in the text.
\end{abstract}

\section{The main theorem and its context}\label{sec:intro}

For a tensor $T\in V_1\otimes\cdots\otimes V_k$, let $\R(T)$ be
its rank with respect to the indicated $k$ factors. We use the grouped
tensor power
\[
 T^{\boxtimes m}\in V_1^{\otimes m}\otimes\cdots\otimes V_k^{\otimes m}
\]
and the asymptotic rank
\begin{equation}\label{eq:asymrank}
 \AR(T)=\lim_{m\to\infty}\R(T^{\boxtimes m})^{1/m},
 \qquad \AR(0)=0.
\end{equation}
For a nonzero tensor the limit exists by submultiplicativity of rank.
Throughout the paper the ground field is $\C$.

\begin{theorem}\label{thm:main}
In the $54$-dimensional vector space
\[
 V=\C^3\otimes\C^3\otimes\C^3\otimes\C^2,
 \qquad a=\sqrt[3]{62},
\]
the asymptotic-rank sublevel
\begin{equation}\label{eq:Y}
 Y=\{T\in V:\AR(T)\leq a\}
\end{equation}
is reducible in the Zariski topology. In particular, the assertion that
all asymptotic-rank sublevels are irreducible, for every algebraically
closed field, every tensor order $k\geq3$, every tuple of positive
factor dimensions, and every real threshold $a\geq0$, is false.
\end{theorem}

The counterexample has order four. The theorem does not settle the
question restricted to order three. Known closedness of asymptotic-rank
sublevels identifies $Y$ as a reducible algebraic variety; the proper
relatively closed cover constructed below proves reducibility directly.

\subsection{The question and the classical ingredients}\label{sec:history}
Christandl, Hoeberechts, Nieuwboer, Vrana and Zuiddam proved that, for
every field, every fixed tensor format of order at least three, and
every real threshold, the corresponding asymptotic-rank sublevel is
Zariski closed \cite[Theorem~1.2]{CHNVZ2024}.
They then asked whether these sublevels are irreducible over
algebraically closed fields \cite[Section~5]{CHNVZ2024}.
The source explicitly distinguishes the familiar irreducible
determinantal varieties of matrices from the open question for
orders $k\geq3$.  The public catalogue records the latter with its
universal quantifiers \cite{CatalogueIrreducibility}.
The present example answers that all-orders question negatively in
order four.  It leaves the separately restricted three-mode question
open.

The asymptotic-rank estimate uses Strassen's asymptotic spectrum and
his entropy formula for tight three-mode supports
\cite{Strassen1988,Strassen1991}.  We use the formulations in
Christandl, Vrana and Zuiddam
\cite[Theorem~1.1, Proposition~1.6, Definition~4.1,
Theorem~4.4 and Corollary~4.5]{CVZ2018}.
Spectral points turn direct sums and grouped products into addition
and multiplication.  The tight-support theorem supplies a lower
bound on an auxiliary tensor's asymptotic subrank, which can then
be subtracted from a direct-sum spectral inequality.  Both classical
theorems are applied to three-mode tensors.  Appending the fourth
factor in the geometric example uses elementary rank and flattening
arguments.

The six-permutation tensor $P$ defined in~\eqref{eq:Pxy} is the variant denoted $\mathrm{cw}'_2$
in the Coppersmith--Winograd literature
\cite{CoppersmithWinograd1990,AlmanLi2026}.
Alman and Li developed quantitative asymptotic-rank speedups using
degenerations with nontrivial direct sums on both sides
\cite[Sections~4 and~7.1]{AlmanLi2026}.  Their Corollary~7.2 in the
pinned arXiv version, equivalently Corollary~34 in the conference
version, applies to this exact tensor and already gives
$\widetilde R(P)\leq\gamma_2'\approx3.9309$ (their Table~1), which is
below $\sqrt[3]{62}$; that bound alone places $x$ in the sublevel.  They also reproduce the classical
$(q+2)$-term Laurent degeneration of the small
Coppersmith--Winograd tensor, crediting Coppersmith and Winograd.

Section~\ref{sec:exchange} gives a self-contained derivation by a direct
finite Fourier exchange, with the exact estimate
\[
 \widetilde R(P)^3\leq74-3\sqrt[3]{81}<62.
\]
The exchange is related to the direct-sum speedup method: for example,
Alman and Li's Proposition~7.1 in the arXiv version
(Proposition~32 in the conference version) exchanges auxiliary matrix
tensors with dimensions $18$ and $8$ when $q=2$ and $n=3$.
The identity proved below instead displays the auxiliary paired
Coppersmith--Winograd tensors and all Fourier coefficients explicitly.
We use the estimate in an ambient open-set argument to produce two
proper relatively closed subsets covering the full asymptotic-rank
sublevel.  The comparisons here are confined to the cited sources;
no priority claim is made for the tensor constructions or rank
estimates.


\subsection{The two witnesses and the structure of the proof}

Let $e_0,e_1,e_2$ be the standard basis of $\C^3$, and let $d_0,d_1$
be the standard basis of $\C^2$. Define
\begin{equation}\label{eq:Pxy}
 \begin{split}
 P&=\sum_{(i,j,k)\text{ a permutation of }(0,1,2)}
             e_i\otimes e_j\otimes e_k,\qquad x=P\otimes d_0,\\
 y&=e_0\otimes e_0\otimes e_0\otimes d_0
       +e_1\otimes e_1\otimes e_1\otimes d_1.
 \end{split}
\end{equation}
Section~\ref{sec:exchange} proves
\begin{equation}\label{eq:strict-cost}
 \AR(P)^3\leq 74-3\sqrt[3]{81}<62,
\end{equation}
so $x\in Y$. The displayed decomposition and a flattening give
$\AR(y)=2$, hence $y\in Y$ as well.

Let $X\subseteq V$ be the closed locus of fourth-mode flattening rank
at most one. We construct a polynomial open set $U\subseteq V$ such
that $x\in U$ and $Y\cap U\subseteq X$. Its defining condition is
the persistence of a one-dimensional kernel of a matrix of appended
minors. Since $y\notin X$, the identity
\begin{equation}\label{eq:cover-intro}
 Y=(Y\cap X)\ \cup\ \bigl(Y\cap(V\setminus U)\bigr)
\end{equation}
is a cover by two proper relatively closed subsets. In particular,
the open set is taken in $V$, without imposing a fixed linear slice
on the sublevel problem.

\section{The asymptotic spectrum and tight supports}\label{sec:classical}

\subsection{Restrictions and the full three-mode spectrum}

A \emph{restriction} $S\leq T$ means that $S$ is obtained from $T$
by a linear map on each factor. The direct sum $S\oplus T$ uses the
direct sum of the respective factor spaces; the product $S\boxtimes T$
groups corresponding factors. These operations make the equivalence
classes under mutual restriction of all complex three-mode tensors
into a preordered semiring. Put
\[
 \Delta_r=\sum_{j=1}^r e_j\otimes e_j\otimes e_j.
\]
A spectral point on this semiring is a map $\phi$ to the nonnegative
real numbers that is additive under $\oplus$, multiplicative under
$\boxtimes$, normalized by $\phi(\Delta_r)=r$, and monotone under
restriction. The spectrum in the next theorem is the spectrum of
\emph{all} complex three-mode tensors.

\begin{theorem}[Strassen's spectral characterization]\label{thm:spectrum}
For every complex three-mode tensor $T$,
\begin{equation}\label{eq:spectrum}
 \AR(T)=\max_{\phi}\phi(T).
\end{equation}
\end{theorem}

We use the formulation and tensor-semiring application in
\cite[Theorem~1.1 and Proposition~1.6]{CVZ2018}.
For later use, the subrank $\Q(T)$ is the largest integer $r$ for
which $\Delta_r\leq T$, and
\[
 \AQ(T)=\lim_{m\to\infty}\Q(T^{\boxtimes m})^{1/m}.
\]
The limit follows from supermultiplicativity for nonzero tensors.
Every spectral point satisfies
\begin{equation}\label{eq:spectral-bounds}
 \AQ(T)\leq\phi(T)\leq\R(T).
\end{equation}
Indeed, a rank decomposition realizes $T$ as a restriction of
$\Delta_{\R(T)}$. Conversely,
$\Q(T^{\boxtimes m})\leq\phi(T)^m$ for every $m$; taking $m$th
roots gives the lower bound. Invertible changes of basis and nonzero
scalar multiplication preserve all spectral values, ranks and
asymptotic ranks, since they give restrictions in both directions.

\subsection{Polynomial degenerations and interpolation}

We will only use explicitly given degenerations. The following
elementary argument explains why each spectral point is monotone
under them.

\begin{lemma}\label{lem:interpolation}
Suppose that for every nonzero complex $t$, linear maps on the three
factors of $K$ produce the tensor
\[
 F(t)=T+tT_1+\cdots+t^dT_d.
\]
The maps may have Laurent-polynomial entries. Then
$\phi(T)\leq\phi(K)$ for every spectral point $\phi$.
\end{lemma}

\begin{proof}
For each positive integer $m$, the polynomial
$F(t)^{\boxtimes m}$ has degree at most $dm$ and constant coefficient
$T^{\boxtimes m}$. Choose $dm+1$ distinct nonzero complex values of
$t$. Polynomial interpolation expresses that constant coefficient
as a linear combination of the values at those points. Each value
is a restriction of $K^{\boxtimes m}$. A linear combination of
$dm+1$ such restrictions is a restriction of their direct sum:
on each direct-sum block use its three maps, absorbing the scalar
coefficient into the first map. Consequently
\[
 \phi(T)^m\leq (dm+1)\phi(K)^m.
\]
Taking $m$th roots and letting $m$ tend to infinity proves the claim.
If $\phi(K)=0$, the same displayed inequality already implies
$\phi(T)=0$. Laurent entries cause no difficulty because the
evaluations avoid zero and the \emph{combined output} is polynomial.
\end{proof}

\subsection{The three-mode tight-support theorem}

Fix bases indexed by finite sets $I_1,I_2,I_3$. A support
$\Phi\subseteq I_1\times I_2\times I_3$ is \emph{tight} if there
are injective maps $\alpha_j:I_j\to\mathbb Z$ such that
\begin{equation}\label{eq:tight}
 \alpha_1(i_1)+\alpha_2(i_2)+\alpha_3(i_3)=0
 \quad\text{for every }(i_1,i_2,i_3)\in\Phi.
\end{equation}
For a probability distribution $p$ on $\Phi$, write $p_j$ for its
marginal on $I_j$, and use natural logarithms in the entropy
$H(p_j)=-\sum_i p_j(i)\log p_j(i)$, with $0\log0=0$.

\begin{theorem}[Strassen's tight-support theorem]\label{thm:tight}
Let $T$ be a complex three-mode tensor whose support in some bases is
tight. For every probability distribution $p$ on that support,
\begin{equation}\label{eq:tight-bound}
 \AQ(T)\geq \min_{1\leq j\leq3}\exp H(p_j).
\end{equation}
\end{theorem}

This is the needed consequence of
\cite[Definition~4.1, Theorem~4.4 and Corollary~4.5]{CVZ2018}.
The hypothesis concerns the support and injective integer weights;
there is no extra equality or positivity hypothesis on its nonzero
coefficients. The theorem is a three-mode result. We apply it to a
three-mode auxiliary tensor, before adjoining the fourth factor in
Section~\ref{sec:cover}.

\section{A strict asymptotic-rank bound for the permutation tensor}
\label{sec:exchange}

\subsection{Paired Coppersmith--Winograd tensors}\label{sec:cw}

For an even positive integer $q$, let $J$ be a set of size $q$
with a fixed-point-free involution $r\mapsto r'$. In a space with
basis $z_0$ and $z_r$ for $r\in J$, define
\begin{equation}\label{eq:paired-cw}
 C_q=\sum_{r\in J}
 \bigl(z_0\otimes z_r\otimes z_{r'}
       +z_r\otimes z_0\otimes z_{r'}
       +z_r\otimes z_{r'}\otimes z_0\bigr).
\end{equation}
Different labelings of the pairs give equivalent tensors.
The usual small Coppersmith--Winograd tensor is
\begin{equation}\label{eq:small-cw}
 \operatorname{cw}_q=\sum_{j=1}^q
 \bigl(e_0\otimes e_j\otimes e_j
       +e_j\otimes e_0\otimes e_j
       +e_j\otimes e_j\otimes e_0\bigr).
\end{equation}

\begin{lemma}\label{lem:cw-basis}
Over $\C$, the tensors $C_q$ and $\operatorname{cw}_q$ are
equivalent under invertible changes of basis. Every spectral point
satisfies $\phi(C_q)\leq q+2$.
\end{lemma}

\begin{proof}
Label the paired vectors $u_h,v_h$, for $1\leq h\leq q/2$.
In each of the three modes apply the map
\[
 z_0\longmapsto\tfrac12e_0,\qquad
 u_h\longmapsto e_{2h-1}+\mathrm i e_{2h},\qquad
 v_h\longmapsto e_{2h-1}-\mathrm i e_{2h}.
\]
The two-by-two block on a pair has determinant $-2\mathrm i$,
and the zero coordinate is scaled by $1/2$, so the map is invertible.
For each pair,
\[
 u_h\otimes v_h+v_h\otimes u_h
 \longmapsto
 2(e_{2h-1}\otimes e_{2h-1}+e_{2h}\otimes e_{2h}).
\]
The factor two is canceled by the zero-coordinate scaling, in each
of its three possible positions. Thus the image is exactly
$\operatorname{cw}_q$.

For the bound, set $s=\sum_{j=1}^q e_j$. For $t\ne0$ the expression
\begin{equation}\label{eq:cw-degeneration}
 t^{-2}\sum_{j=1}^q(e_0+te_j)^{\otimes3}
 -t^{-3}(e_0+t^2s)^{\otimes3}
 +(t^{-3}-qt^{-2})e_0^{\otimes3}
\end{equation}
is a sum of at most $q+2$ rank-one tensors. Expanding it gives
\begin{equation}\label{eq:cw-polynomial}
 \operatorname{cw}_q
 +t\left(\sum_{j=1}^q e_j^{\otimes3}
       -e_0\otimes s\otimes s-s\otimes e_0\otimes s
       -s\otimes s\otimes e_0\right)-t^3s^{\otimes3}.
\end{equation}
The coefficients of $e_0^{\otimes3}$ cancel, and the degree $-1$
terms with two zero coordinates cancel. The remaining degree-zero
terms are exactly~\eqref{eq:small-cw}. Applying
Lemma~\ref{lem:interpolation} with source $\Delta_{q+2}$ shows
$\phi(\operatorname{cw}_q)\leq q+2$. Basis invariance proves the
same bound for $C_q$.
\end{proof}

\subsection{The group tensor and its Fourier decomposition}

Identify $\F_2^2$ with the symbols $0,1,2,3$ by their two-bit
expansions, so addition is bitwise exclusive-or. Set
\begin{equation}\label{eq:sets}
 \begin{split}
 G&=(\F_2^2)^3,\qquad S=\{1,2,3\}^3,\qquad b=(1,1,1),\\
 I&=S\cap(S+b),\qquad E=G\setminus\bigl(S\cup(S+b)\bigr).
 \end{split}
\end{equation}
Here $I=\{2,3\}^3$. Since $|G|=64$ and $|S|=27$,
inclusion-exclusion gives
\begin{equation}\label{eq:set-sizes}
 |I|=8,\qquad |E|=64-(27+27-8)=18.
\end{equation}
Both $I$ and $E$ are invariant under $r\mapsto r+b$.
This is a fixed-point-free involution because $b\ne0$.
Also $0\notin S\cup E$: indeed $b\in S$ implies $0\in S+b$.

Let $f_g$, for $g\in G$, denote a basis and define the group tensor
\begin{equation}\label{eq:group-tensor}
 K_G=\sum_{u+v+w=0}f_u\otimes f_v\otimes f_w.
\end{equation}
Its Walsh decomposition is
\begin{equation}\label{eq:walsh}
 K_G=\frac1{64}\sum_{\chi\in\widehat G}
       \left(\sum_{g\in G}\chi(g)f_g\right)^{\otimes3},
\end{equation}
where $\widehat G$ is the set of the $64$ characters with values
in $\{1,-1\}$. For a coefficient indexed by $(u,v,w)$ the right
side is $64^{-1}\sum_\chi\chi(u+v+w)$. This equals one if the
sum is zero. Otherwise, choosing a character that is negative on
that nonzero group element pairs the terms with their negatives.
This proves~\eqref{eq:walsh} and $\R(K_G)\leq64$.

\subsection{The finite exchange}\label{sec:finite-exchange}

In each target factor use the disjoint basis vectors
\[
 \{x_g:g\in S\}\ \sqcup\ \{z_0\}\ \sqcup\ \{z_g:g\in E\}.
\]
Apply in all three modes the Laurent map
\begin{equation}\label{eq:laurent-map}
 f_g\longmapsto
 \boldsymbol1_{g\in S}x_g
 +t^{-2}\boldsymbol1_{g=b}z_0
 +t\boldsymbol1_{g\in E}z_g.
\end{equation}
In particular, $f_b$ has both an $x_b$ component and a $z_0$
component; the formula is a sum of components, not a partition of
the source basis. The image of $f_0$ is zero.

\begin{lemma}\label{lem:exchange}
There are Laurent-polynomial restrictions of $\Delta_{64}\oplus C_8$
whose combined output is a polynomial in $t$ of degree at most three
with constant coefficient
\begin{equation}\label{eq:exchange-target}
 P^{\boxtimes3}\oplus C_{18}.
\end{equation}
Consequently every spectral point satisfies
\begin{equation}\label{eq:exchange-spectral}
 \phi(P)^3+\phi(C_{18})\leq64+\phi(C_8).
\end{equation}
\end{lemma}

\begin{proof}
First apply~\eqref{eq:laurent-map} to $K_G$. A term using only
$x$ components has degree zero. For each of the three coordinates
of $G$, three elements of $\{1,2,3\}$ sum to zero exactly when
they are the three distinct nonzero elements of $\F_2^2$.
Identify $1,2,3$ with the basis labels $0,1,2$ of $P$.
The $x$-only part is then $P^{\boxtimes3}$, with $6^3=216$
nonzero coefficients, all equal to one.

Consider a term with exactly one $z_0$ component. Its corresponding
group index is $b$, and the other two group indices must sum to $b$.
There are three possibilities for their target types:
\begin{enumerate}
\item Two $x$ components have degree $-2$. Their indices lie in
$I$, paired by $r\mapsto r+b$. Taking all three positions of $z_0$
gives exactly $t^{-2}C_8(z_0,x_i:i\in I)$.
\item One $x$ component and one $z_r$, for $r\in E$, would have
degree $-1$, but cannot occur. If the first index is in $S$, the
second must be in $S+b$, which is disjoint from $E$.
\item Two $E$ components have degree zero. Their indices are the
$18$ elements of $E$, paired by the same involution. These terms
give exactly $C_{18}$, with $3\cdot18=54$ support coefficients.
\end{enumerate}
Two $z_0$ components would require the third group index to be zero,
whose image is zero. Three would require $b+b+b=0$, contrary to
$b\ne0$. Finally, terms with no $z_0$ and at least one $E$ component
have positive degree, at most three. This exhausts every possible
term, including the two components of the image of $f_b$.

Now adjoin a disjoint copy of the paired tensor $C_8$ to the source.
Label its nonzero basis vectors by $I$ with the pairing $i\mapsto i+b$.
In each mode send its zero-index basis vector to $-t^{-2}z_0$ and its
other basis vectors to $x_i$. Every term of that tensor has exactly
one zero-index basis vector, so its image is precisely
$-t^{-2}C_8(z_0,x_i:i\in I)$. This cancels all negative degrees.
The resulting output is polynomial of degree at most three, with
the constant coefficient in~\eqref{eq:exchange-target}. Its two
summands occupy disjoint coordinates in every factor.

Equation~\eqref{eq:walsh} realizes $K_G$ as a restriction of
$\Delta_{64}$, so the same construction has source
$\Delta_{64}\oplus C_8$. Lemma~\ref{lem:interpolation}, together
with additivity and multiplicativity of $\phi$, proves
\eqref{eq:exchange-spectral}.
\end{proof}

\subsection{The tight auxiliary tensor and the exact constant}

Choose labels $u_h,v_h$, for $1\leq h\leq9$, for the pairs of
$C_{18}$. Assign weight $0$ to $z_0$, weight $h$ to $u_h$ and
weight $-h$ to $v_h$, in each of the three modes. These weights
are injective within a mode and sum to zero on every support
triple. Thus $C_{18}$ has tight support.

Give each of its $54$ support triples probability $1/54$.
Every marginal has mass $1/3$ at the zero coordinate and mass
$2/54=1/27$ at each of its $18$ other coordinates. Its entropy
therefore satisfies
\begin{equation}\label{eq:entropy}
 \begin{split}
 \exp H
 &=\exp\left(-\tfrac13\log\tfrac13
             -18\cdot\tfrac1{27}\log\tfrac1{27}\right)\\
 &=3^{7/3}=3\sqrt[3]{81}.
 \end{split}
\end{equation}
Theorem~\ref{thm:tight} and~\eqref{eq:spectral-bounds} give
$\phi(C_{18})\geq\AQ(C_{18})\geq3\sqrt[3]{81}$ for every
spectral point. Lemma~\ref{lem:cw-basis} gives $\phi(C_8)\leq10$.
Substituting both inequalities into~\eqref{eq:exchange-spectral}
yields the uniform estimate
\begin{equation}\label{eq:phiP}
 \phi(P)^3\leq64+10-3\sqrt[3]{81}
             =74-3\sqrt[3]{81}.
\end{equation}
Taking the maximum in Theorem~\ref{thm:spectrum} proves the
first inequality of~\eqref{eq:strict-cost}.

All comparisons needed for the threshold are exact. Since
$81>64=4^3$, one has $3\sqrt[3]{81}>12$, so the last expression
in~\eqref{eq:phiP} is strictly less than $62$. Moreover
\begin{equation}\label{eq:threshold}
 27<62<64
 \quad\Longrightarrow\quad
 3<\sqrt[3]{62}<4.
\end{equation}
No decimal estimate or assumption about $\AR(P)$ is used.

\section{Scalar centroids and an ambient open neighborhood}
\label{sec:geometry}

Set $A=B=C=\C^3$ and $D=\C^2$. Every element of $V$ has a
unique expression
\begin{equation}\label{eq:slices}
 W=T\otimes d_0+S\otimes d_1,
 \qquad T,S\in A\otimes B\otimes C.
\end{equation}
The fourth-mode flattening views $W$ as a map
$D^*\to A\otimes B\otimes C$. Its rank is the dimension of
$\operatorname{span}\{T,S\}$. Hence
\begin{equation}\label{eq:X}
 X=\{W:\rank(W_{ABC\mid D})\leq1\}
   =\{T\otimes v:T\in A\otimes B\otimes C,\ v\in D\}
\end{equation}
is closed, defined by the two-by-two minors of a $27$-by-$2$ matrix.

Let $Z\subseteq V$ be the common rank-at-most-three locus of the
three balanced flattenings
\begin{equation}\label{eq:Z}
 AB\mid CD,\qquad AC\mid BD,\qquad BC\mid AD.
\end{equation}
These are $9$-by-$6$ matrices, so their four-by-four minors define
$Z$. Every point of $X$ lies in $Z$: for example the $AB\mid CD$
flattening of $T\otimes(v_0d_0+v_1d_1)$ is
$[v_0M\ v_1M]$, where $M=T_{AB\mid C}$ has three columns.

\subsection{An appended-minor matrix}

Write $M_C(T)$ for the $9$-by-$3$ flattening with row indices in
$A\times B$ and column indices in $C$. Define $M_B(T)$ and
$M_A(T)$ by the other two choices of the singled-out factor.
A tensor $T$ is \emph{concise} if all three matrices have rank three.

For each of these three matrices, append one of the three columns
of the corresponding matrix of $S$, and take the determinant on
each choice of four rows. These expressions are linear in the
$27$ entries of $S$. Collect their coefficients into the matrix
\begin{equation}\label{eq:L}
 L(T):\C^{27}\longrightarrow\C^{1134},
 \qquad 1134=3\cdot3\binom94.
\end{equation}
Every entry of $L(T)$ is a homogeneous cubic polynomial in the
coordinates of $T$. Appending a column of $T$ repeats an existing
column, so
\begin{equation}\label{eq:LT}
 L(T)T=0\quad\text{for every }T.
\end{equation}

\begin{lemma}\label{lem:appended}
If $T$ is concise, then the tensor $W$ in~\eqref{eq:slices}
belongs to $Z$ if and only if $L(T)S=0$.
\end{lemma}

\begin{proof}
Up to column permutations, each balanced flattening of $W$ is
$[M_j(T)\ M_j(S)]$. Since $M_j(T)$ has rank three, that block
matrix has rank at most three exactly when every column of
$M_j(S)$ belongs to the column space of $M_j(T)$. Appending a
single column gives a $9$-by-$4$ matrix. Its four-by-four minors
all vanish exactly when its rank is at most three, which, because
the original three columns are independent, is precisely the
required membership condition. Doing this for all columns and
all three flattenings proves both directions.
\end{proof}

\subsection{The centroid interpretation and the kernel at \texorpdfstring{$P$}{P}}

For a concise tensor $T$, define its centroid to be the vector space
\begin{equation}\label{eq:centroid}
 \begin{split}
 \Cent(T)=\{(F_A,F_B,F_C)\in\End(A)\oplus\End(B)\oplus\End(C):\ &
 (F_A\otimes1\otimes1)T\\
 &=(1\otimes F_B\otimes1)T
 =(1\otimes1\otimes F_C)T\}.
 \end{split}
\end{equation}
Taking the common value identifies this vector space with
$\ker L(T)$. To see this, concision makes each map
$F_j\mapsto F_j\cdot T$ injective. Its image consists precisely
of those $S$ for which the columns of $M_j(S)$ belong to the
column space of $M_j(T)$: a matrix relation
$M_j(S)=M_j(T)H$ corresponds to action by the transpose of $H$
on the singled-out factor. The three simultaneous membership
conditions are exactly those defining $\ker L(T)$.
Under this identification, the common scalar triples
$(\lambda1_A,\lambda1_B,\lambda1_C)$ correspond to $\lambda T$.

\begin{lemma}\label{lem:scalar}
The tensor $P$ is concise, and
\begin{equation}\label{eq:kernelP}
 \ker L(P)=\C P,\qquad \rank L(P)=26.
\end{equation}
Equivalently, $\Cent(P)$ consists exactly of the common scalar triples.
\end{lemma}

\begin{proof}
The columns of $M_C(P)$ are the three independent vectors
\[
 e_1\otimes e_2+e_2\otimes e_1,\quad
 e_0\otimes e_2+e_2\otimes e_0,\quad
 e_0\otimes e_1+e_1\otimes e_0.
\]
Their span is the space of symmetric three-by-three matrices with
zero diagonal. The same description holds for $M_A(P)$ and
$M_B(P)$ by permutation symmetry, proving concision.

If $L(P)S=0$, the column-space characterization in the preceding
lemma says, for each fixed $k$, that the array $(S_{ijk})_{i,j}$
is symmetric with zero diagonal. The other two flattenings give
the same condition for the other pairs of indices. Thus the
coefficients of $S$ are unchanged by each transposition of the
three indices and vanish whenever any two indices coincide.
The only possible nonzero coefficients are the six permutations
of $(0,1,2)$; transpositions act transitively on these six, so all
six coefficients agree. Therefore $S$ is a scalar multiple of $P$.
Conversely,~\eqref{eq:LT} shows that every such multiple lies in
the kernel. Rank-nullity gives $27-1=26$, and the centroid
description follows from the preceding identification.
\end{proof}

\subsection{Persistence in the full ambient space}

Define $U\subseteq V$ by the following conditions on the first
slice in~\eqref{eq:slices}:
\begin{equation}\label{eq:U}
 U=\{T\otimes d_0+S\otimes d_1:
       \rank M_A(T)=\rank M_B(T)=\rank M_C(T)=3,
       \ \rank L(T)\geq26\}.
\end{equation}
Each condition is polynomial open: maximal flattening rank is
the nonvanishing of at least one three-by-three minor, and
$\rank L(T)\geq26$ is the nonvanishing of at least one
$26$-by-$26$ minor. The slice projection $W\mapsto T$ is linear,
so these conditions indeed define an open subset of all of $V$.
Lemma~\ref{lem:scalar} shows that $x=P\otimes d_0\in U$.

\begin{proposition}\label{prop:open}
The open set $U$ satisfies $Z\cap U\subseteq X$.
\end{proposition}

\begin{proof}
Take $W=T\otimes d_0+S\otimes d_1\in Z\cap U$.
Concision implies $T\ne0$. By~\eqref{eq:LT}, the kernel of
$L(T)$ contains this nonzero vector, so its rank is at most $26$.
The defining inequality of $U$ forces equality. Thus
$\ker L(T)=\C T$. Since $W\in Z$, Lemma~\ref{lem:appended}
gives $L(T)S=0$, and therefore $S=\lambda T$ for a scalar
$\lambda$. It follows that
\[
 W=T\otimes(d_0+\lambda d_1)\in X.
\]
\end{proof}

This is the only persistence statement needed: one nonzero minor
at $P$, together with the universal kernel vector $T$, keeps the
kernel one-dimensional on an open neighborhood. No assertion
about arbitrary tensors with deficient single-mode ranks is used.

\section{The proper relatively closed cover}\label{sec:cover}

We first record why all relevant matrix-rank constraints follow
from asymptotic tensor rank, and why adding a fixed factor does
not change the rank of a tensor.

\begin{lemma}\label{lem:flatten}
For any tensor $T$ of any order, the rank of every matrix
flattening of $T$ is at most $\AR(T)$.
\end{lemma}

\begin{proof}
Let a flattening have rank $r$. The corresponding flattening of
$T^{\boxtimes m}$ is, up to permutations of rows and columns,
the $m$th matrix Kronecker power of that flattening and has
rank $r^m$. Every rank-one tensor becomes a matrix of rank at
most one under a flattening, so $r^m\leq\R(T^{\boxtimes m})$.
Taking $m$th roots and the limit proves the claim, including the
zero tensor.
\end{proof}

\begin{lemma}\label{lem:append}
If $v$ is a fixed nonzero vector in an additional factor, then
$\R(T\otimes v)=\R(T)$ and $\AR(T\otimes v)=\AR(T)$,
where the two ranks use their respective numbers of factors.
\end{lemma}

\begin{proof}
A rank decomposition of $T$ gives one of $T\otimes v$ with
the same number of terms. Choose a linear functional $\ell$
with $\ell(v)=1$. Applying it to the last factor of any
decomposition of $T\otimes v$ gives a decomposition of $T$,
with no more terms. These are the two rank inequalities.
Apply the same argument to each power, with added vector
$v^{\otimes m}$, and take $m$th roots and limits.
\end{proof}

\begin{proof}[Proof of Theorem~\ref{thm:main}]
For any $W\in Y$, Lemma~\ref{lem:flatten} bounds each balanced
flattening rank by $\AR(W)\leq a<4$. Matrix rank is integral,
so all three ranks are at most three. Therefore $Y\subseteq Z$.

Equation~\eqref{eq:strict-cost} and Lemma~\ref{lem:append} imply
$\AR(x)=\AR(P)<a$, so $x\in Y$. The tensor $y$ in
\eqref{eq:Pxy} has rank at most two by its displayed decomposition.
Its fourth-mode flattening has two independent columns, so
Lemma~\ref{lem:flatten} gives $\AR(y)\geq2$. Since
$\AR(y)\leq\R(y)\leq2$, equality holds. In particular
$y\in Y$, because $2<a$, and $y\notin X$.

The sets
\[
 Y_1=Y\cap X,\qquad Y_2=Y\cap(V\setminus U)
\]
are relatively Zariski closed in $Y$, since $X$ and $V\setminus U$
are closed in $V$. They cover $Y$: if a point of $Y$ lies in $U$,
then it lies in $Z\cap U\subseteq X$ by
Proposition~\ref{prop:open}; otherwise it belongs to $Y_2$.
Both pieces are proper. The point $y$ belongs to $Y\setminus Y_1$,
and the point $x$ belongs to $Y\setminus Y_2$ because $x\in U$.
Thus $Y$ is reducible as a topological subspace of $V$.
Its known Zariski closedness \cite[Theorem~1.2]{CHNVZ2024}
then gives the stated interpretation as a reducible algebraic
variety.
\end{proof}

The use of four factors occurs in the pair of slices and their
balanced flattenings. The spectral and tight-support estimates
are applied exclusively to three-mode tensors. The proper cover
does not imply that either displayed piece is irreducible or
that there are exactly two irreducible components.

\section{Component growth and a boundary under grouping}
\label{sec:components}

For complex vector spaces $V_1,\ldots,V_k$ and a real $a\geq0$, write
\[
 Z_a(V_1,\ldots,V_k)
   =\{T\in V_1\otimes\cdots\otimes V_k:\AR(T)\leq a\}.
\]
These are the full sublevels, which are Zariski closed by
\cite[Theorem~1.2]{CHNVZ2024}. Ranks and asymptotic ranks always refer
to the displayed factors. Grouping a collection of factors for a matrix
flattening does not change this convention.

\subsection{The local rank-one chart in general dimension}
\label{sec:components-chart}

Let $A,B,C$ have the same dimension $n\geq3$, put $W=A\otimes B\otimes C$,
and let $D$ have dimension $m\geq2$. For $T\in W$, set
\[
 E_{AB}(T)=\im(C^*\longrightarrow A\otimes B),
\]
and define $E_{AC}(T)$ and $E_{BC}(T)$ in the same way. Using the
natural permutations of tensor factors, define
\begin{equation}\label{eq:components-K}
 K(T)=(E_{AB}(T)\otimes C)\cap(E_{AC}(T)\otimes B)
                    \cap(A\otimes E_{BC}(T))\subseteq W.
\end{equation}
Thus $T\in K(T)$, and $T$ is concise precisely when the three
pair-image spaces have dimension $n$.

The appended-minor construction of Section~\ref{sec:geometry} extends
without change: to each $n^2$-by-$n$ flattening of $T$, append each
column of the corresponding flattening of a variable $S$, and take
all $(n+1)$-row determinants. Their coefficients in $S$ form a
polynomial matrix $L_n(T)$ with $n^3$ columns. Repeated columns give
$L_n(T)T=0$ for every $T$, and on the concise locus one has
\[
 \ker L_n(T)=K(T).
\]
Indeed, a full-rank $n$-column matrix acquires no additional rank
on appending a column exactly when that column lies in its image.
Consequently
\begin{equation}\label{eq:components-Omega}
 \Omega_n=\{T\in W:T\text{ is concise},\ \rank L_n(T)\geq n^3-1\}
\end{equation}
is Zariski open, and $K(T)=\C T$ for every $T\in\Omega_n$.

Choose a basis $d_0,\ldots,d_{m-1}$ of $D$ and write
$F=\sum_{j=0}^{m-1}T_j\otimes d_j$. Define
\[
 \begin{split}
 \mathcal F_n={}&\{F:\rank F_{AB\mid CD},\ \rank F_{AC\mid BD},
                          \ \rank F_{BC\mid AD}\leq n\},\\
 X={}&\{F:\rank F_{ABC\mid D}\leq1\},\qquad
 U=\{F:T_0\in\Omega_n\}.
 \end{split}
\]
Here the first line requires each of the three ranks to be at most $n$.

\begin{proposition}\label{prop:components-chart}
The full-ambient open set $U$ satisfies
\begin{equation}\label{eq:components-chart-equality}
 \mathcal F_n\cap U=X\cap U.
\end{equation}
If $\Omega_n\ne\varnothing$, then $X$ is an irreducible component of
$\mathcal F_n$ of dimension $n^3+m-1$. In the four-factor rank convention
$(A,B,C,D)$, for every $a\geq0$ one has
\begin{equation}\label{eq:components-sublevel-X}
 Z_a(A,B,C,D)\cap X
   =\{T\otimes v:T\in Z_a(A,B,C),\ v\in D\}.
\end{equation}
For $0\leq a<n+1$, its intersection with $U$ is isomorphic to
\begin{equation}\label{eq:components-sublevel-chart}
 \bigl(Z_a(A,B,C)\cap\Omega_n\bigr)\times\mathbb A^{m-1}.
\end{equation}
\end{proposition}

\begin{proof}
Suppose $F\in U$. The $AB\mid CD$ flattening concatenates the columns
of all the $T_j$, while the first slice already has an $n$-dimensional
$AB$-image. This flattening has rank at most $n$ if and only if
$E_{AB}(T_j)\subseteq E_{AB}(T_0)$ for every $j$. The other two balanced
flattenings give the corresponding conditions on $AC$ and $BC$.
Therefore $F\in\mathcal F_n$ precisely when every $T_j$ belongs to
$K(T_0)=\C T_0$, which is precisely $F\in X$. Conversely, every
$T\otimes v$ has each balanced rank at most $n$, so $X\subseteq\mathcal F_n$.
This proves~\eqref{eq:components-chart-equality}.

The locus $X$ is the rank-at-most-one variety of $n^3$-by-$m$ matrices.
It is closed by its two-by-two minors and irreducible as the image of
$W\times D$ under $(T,v)\mapsto T\otimes v$. Its dimension is
$n^3+m-1$: on a chart with a chosen nonzero matrix entry, a rank-one
matrix is determined by its corresponding row and column, with their
common entry counted once. If an irreducible closed subset $C$ of
$\mathcal F_n$ contains $X$, then $C\cap U$ is nonempty when
$\Omega_n\ne\varnothing$. It is dense in $C$, and
\eqref{eq:components-chart-equality} places it in the closed subset
$X$. Hence $C=X$, proving the component assertion.

For $v\ne0$, Lemma~\ref{lem:append} gives
$\AR(T\otimes v)=\AR(T)$; the zero tensor is covered by taking $T=0$.
This proves~\eqref{eq:components-sublevel-X}. If $a<n+1$,
Lemma~\ref{lem:flatten} and integrality give
$Z_a(A,B,C,D)\subseteq\mathcal F_n$. The isomorphism in
\eqref{eq:components-sublevel-chart} is
\[
 (T,\lambda_1,\ldots,\lambda_{m-1})
   \longmapsto T\otimes\left(d_0+\sum_{j=1}^{m-1}\lambda_jd_j\right).
\]
Its inverse reads the first slice $T\ne0$ and the unique scalar multiples
of it in the other slices. Those scalars are regular on each chart on
which a selected coordinate of $T$ is nonzero, and these expressions
agree on overlaps. This also proves the claimed isomorphism globally
on the indicated open subset.
\end{proof}

The component assertion concerns the determinantal locus $\mathcal F_n$;
it does not assert that the whole of $X$ lies in an asymptotic-rank
sublevel. Equation~\eqref{eq:components-sublevel-X} retains the actual
three-mode sublevel. The exact sublevel identifications above use four
factors, with $D$ a single factor. When $D$ abbreviates several original
factors below, only the flattening statement
\eqref{eq:components-chart-equality} is used.

\subsection{Scalar-kernel cores in every dimension}
\label{sec:general-cores}

In this subsection $A=B=C=\C^n$, and $K(T)$ is the intersection
of lifted two-mode images defined above. We write $\underline{\R}(T)$
for border rank: the least integer $r$ for which $T$ belongs to the
Zariski closure of the rank-at-most-$r$ locus.

For an integer $n\geq3$ and $s\in\C\setminus\{0\}$, put
\begin{equation}\label{eq:general-core}
 w=\sum_{i=1}^n e_i,
 \qquad T_{n,s}=\sum_{i=1}^n e_i\otimes e_i\otimes e_i
                    +s\,w\otimes w\otimes w.
\end{equation}

\begin{lemma}\label{lem:general-core}
For every $n\geq3$ and every nonzero complex $s$, the tensor
$T_{n,s}$ is concise and $K(T_{n,s})=\C T_{n,s}$.
\end{lemma}

\begin{proof}
Its slices in any mode are $E_{ii}+sJ$, where $J$ is the all-ones
matrix. A linear relation among these slices has the form
\[
 \operatorname{diag}(x_1,\ldots,x_n)
                      +s\Bigl(\sum_i x_i\Bigr)J=0.
\]
An off-diagonal entry gives $\sum_i x_i=0$, and the diagonal entries
then give $x_i=0$ for every $i$. This proves concision in each mode.

Let $S\in K(T_{n,s})$. Every matrix in each two-mode image is
symmetric, so the three lifted-image conditions make $S$ invariant
under all permutations of its indices. In the slice with third
index $k$, write $q_k$ for the common off-diagonal entry. Expressing
this slice in the basis $E_{ii}+sJ$ gives
\[
 q_k=s\left(\sum_i S_{iik}-nq_k\right).
\]
For distinct $j,k$, choose $i$ different from both. Symmetry gives
$S_{ijk}=S_{ikj}$, hence $q_k=q_j$. Let their common value be $q$.
For $i\ne k$, symmetry also gives $S_{iik}=S_{iki}=q$.
Writing $b_k=S_{kkk}$, the displayed relation becomes
$q=s(b_k-q)$, so $b_k=(1+1/s)q$. Thus the coefficients of $S$ are
$q$ off the full diagonal and $(1+1/s)q$ on it, which means
$S=(q/s)T_{n,s}$. The reverse inclusion follows from
$T_{n,s}\in K(T_{n,s})$.
\end{proof}

\begin{lemma}\label{lem:general-core-rank}
For every $n\geq3$ and every nonzero complex $s$,
\begin{equation}\label{eq:general-core-ranks}
 \R(T_{n,s})=\underline{\R}(T_{n,s})=n+1.
\end{equation}
\end{lemma}

\begin{proof}
Equation~\eqref{eq:general-core} gives the upper bound $n+1$ for
both ranks. A concise tensor of rank at most $n$ has rank exactly
$n$, since each single-mode flattening has rank $n$. In an
$n$-term decomposition, its factors span each of the three
$n$-dimensional modes and hence form bases. It is therefore an
invertible local transform of
\[
 D_n=\sum_{i=1}^n e_i\otimes e_i\otimes e_i.
\]
For $D_n$, the lifted-image conditions force all three indices to
coincide, so $K(D_n)$ is the $n$-dimensional span of its diagonal
summands. Invertible local maps carry $K(T)$ to the corresponding
kernel for the transformed tensor. Consequently every concise
tensor of rank at most $n$ has kernel dimension $n$.

On the concise open subset, the appended-minor construction of
Section~\ref{sec:geometry} works with $n^2$-by-$n$ flattenings:
append each column of a variable tensor and impose every
$(n+1)$-row determinant. The resulting polynomial matrix $L_n(T)$
has $n^3$ columns and satisfies $\ker L_n(T)=K(T)$ there.
The condition $\dim K(T)\geq n$ is therefore relatively closed,
being the condition $\rank L_n(T)\leq n^3-n$.
Rank-at-most-$n$ tensors remain dense in their border-rank closure
after restriction to this open subset. Every concise tensor of
border rank at most $n$ must consequently satisfy
$\dim K(T)\geq n$. Lemma~\ref{lem:general-core} excludes this
possibility for $T_{n,s}$, whose kernel has dimension one.
Thus $\underline{\R}(T_{n,s})>n$, proving both equalities.
\end{proof}

\begin{proposition}\label{prop:components-core}
For every integer $n\geq3$, define
\begin{equation}\label{eq:components-Bn}
 B_n=\left((n+1)^n+n^{2n/3}
            -\bigl((n+1)^n-n^n\bigr)^{2/3}\right)^{1/n},
\end{equation}
where the roots are positive real roots. Then $n<B_n<n+1$, and
for every $s\in\C\setminus\{0\}$ the tensor $T_{n,s}$ is concise,
with
\[
 K(T_{n,s})=\C T_{n,s},\qquad
 \R(T_{n,s})=\underline{\R}(T_{n,s})=n+1,\qquad
 \AR(T_{n,s})\leq B_n.
\]
\end{proposition}

\begin{proof}
The concision, kernel and exact-rank assertions are
Lemmas~\ref{lem:general-core} and~\ref{lem:general-core-rank}.
For the asymptotic bound we use the automatic three-direction
form of the speedup theorem of Alman and Li \cite[Theorem~6.1]{AlmanLi2026v2}. For a tensor
$Q$ of format $N\times N\times N$ and an integer ordinary-rank
upper bound $R>2N$, its automatic clause gives
\begin{equation}\label{eq:automatic-speedup}
 \AR(Q)\leq R+N^{2/3}-(R-N)^{2/3}<R.
\end{equation}
Here the auxiliary matrix-rank bound is $N$ in each of the three
directions, which is automatic for $N$-by-$N$ matrices. We apply
this theorem only to three-mode tensors.

Take $Q=T_{n,s}^{\boxtimes n}$, $N=n^n$ and $R=(n+1)^n$.
Taking the grouped $n$th power of the explicit $(n+1)$-term
decomposition~\eqref{eq:general-core} gives $\R(Q)\leq R$. The binomial expansion shows
\[
 (1+1/n)^n
   =2+\sum_{j=2}^n\binom nj n^{-j}>2,
\]
so $R>2N$. Grouped powers satisfy
$\AR(T^{\boxtimes n})=\AR(T)^n$, directly from the defining
limit~\eqref{eq:asymrank}. Substitution in
\eqref{eq:automatic-speedup} yields
$\AR(T_{n,s})^n\leq B_n^n$.

For the endpoints put $z=R-N>N>1$. Since
$N^{2/3}<z^{2/3}$, one has $B_n^n<R$. Also
\[
 B_n^n-N=z-z^{2/3}+N^{2/3}>0.
\]
Taking positive $n$th roots gives $n<B_n<n+1$.
\end{proof}

Thus these cores are available throughout the interval
$B_n\leq a<n+1$: the lower endpoint is included because the
asymptotic-rank estimate is non-strict. The upper endpoint is
excluded because the component argument uses integral flattening
ranks bounded above by $a<n+1$. No exact value of
$\AR(T_{n,s})$ is asserted.
For dimension three, the permutation tensor $P$ supplies the
separate interval $\sqrt[3]{62}\leq a<4$: its scalar kernel is
Lemma~\ref{lem:scalar}, and its strict cost and threshold bounds
are~\eqref{eq:strict-cost} and~\eqref{eq:threshold}.

\subsection{A lower bound for the number of components}
\label{sec:components-count}

\begin{theorem}\label{thm:components-count}
Let $n\geq3$ and $k\geq4$ be integers, let $d_1,\ldots,d_k\geq2$, and put
\[
 t_n=\#\{i:d_i\geq n\}.
\]
For every real $a$ satisfying $B_n\leq a<n+1$, with $B_n$ as in
\eqref{eq:components-Bn}, the variety
$Z_a(\C^{d_1},\ldots,\C^{d_k})$ has at least
\begin{equation}\label{eq:components-count}
                         \binom{t_n}{3}+1
\end{equation}
irreducible components. The same bound with $t_3$ holds for every
$\sqrt[3]{62}\leq a<4$.
\end{theorem}

\begin{proof}
First suppose that every $d_i\leq n$, and call a factor active if its
dimension is $n$. For the first interval, choose any nonzero $s$ and
put $Q=T_{n,s}$ from Proposition~\ref{prop:components-core}. For the
second interval put $n=3$ and $Q=P$. Lemma~\ref{lem:scalar} and
\eqref{eq:strict-cost} supply its scalar kernel and cost bound. Thus
in both cases $Q$ is concise, $K(Q)=\C Q$, and
\[
                         \AR(Q)\leq a<n+1.
\]
In particular $a\geq n\geq3$ by Lemma~\ref{lem:flatten}.

For each triple $I$ of active factors, let $x_I$ be $Q$ on those
factors and a fixed basis vector $e_0$ on every other factor.
Repeated application of Lemma~\ref{lem:append} gives
$\AR(x_I)=\AR(Q)$, so $x_I\in Z_a$.
For linear-algebra purposes write
$D=\bigotimes_{j\notin I}\C^{d_j}$ and choose its product basis.
The complementary all-zero coordinate selects the first slice $Q$.
Proposition~\ref{prop:components-chart} supplies a full-ambient open
neighborhood $U_I$ of $x_I$. Every bipartition rank on $Z_a$ is an
integer at most $n$, so
\begin{equation}\label{eq:components-local}
 Z_a\cap U_I\subseteq X_I,
 \qquad X_I=\{F:\rank F_{I\mid I^c}\leq1\}.
\end{equation}
This use of $D$ does not merge any modes in the definition of $Z_a$.

The closed affine variety $Z_a$ has finitely many irreducible components.
Choose a component $C_I$ containing $x_I$. Its nonempty relatively
open subset $C_I\cap U_I$ is dense in $C_I$ and lies in the closed
locus $X_I$. It follows that
\begin{equation}\label{eq:components-containment}
                              C_I\subseteq X_I.
\end{equation}
No uniqueness of a component through $x_I$ is needed.

There is a useful stronger constraint: every single-mode rank outside
$I$ is at most one everywhere on $C_I$. Fix $i\in I$. The locus on
$C_I$ where the $i$th single-mode rank is $n$ is nonempty, since it
contains $x_I$, and is open, since that mode has dimension $n$.
At any point $F$ of this open set, \eqref{eq:components-containment}
gives a nonzero factorization $F=A_I\otimes B_{I^c}$. For $j\notin I$,
the Kronecker-product rank formula gives
\[
 \rank F_{\{i,j\}\mid\{i,j\}^c}
      =r_i(A_I)r_j(B_{I^c})=n\,r_j(B_{I^c})\leq n.
\]
Here $r_j$ denotes the appropriate single-mode flattening rank, and
the last inequality follows from $F\in Z_a$. Thus $r_j(F)=1$ on
that open set. The condition $r_j(F)\leq1$ is closed, so density
extends it to all of $C_I$. If $I\ne J$, choose $j\in J\setminus I$.
Then $x_J$ has $j$th rank $n$, whereas every point of $C_I$ has
$j$th rank at most one. Therefore the chosen components $C_I$ are
pairwise distinct, whether or not the triples overlap.

Every factor has dimension at least two. The tensor
\[
 g=\bigotimes_{i=1}^k e_0+\bigotimes_{i=1}^k e_1
\]
has rank and asymptotic rank two: its displayed decomposition gives
the upper bound, and every nontrivial bipartition has matrix rank two.
Thus $g\in Z_a$. Since $k\geq4$, each $I\mid I^c$ is nontrivial,
so $g\notin X_I$ for every active triple. A component containing $g$
is different from all the $C_I$, proving~\eqref{eq:components-count}
in these dimensions. If there are no active triples, the bound is
one and follows from the same nonempty sublevel.

For arbitrary $d_i\geq2$, choose coordinate projections
\[
 p_i:\C^{d_i}\longrightarrow\C^{\min(d_i,n)}
\]
with coordinate inclusions $j_i$ as sections. The tensor-product maps
$p=\bigotimes_i p_i$ and $j=\bigotimes_i j_i$ are local maps for the
original $k$ factors and satisfy $pj=1$. Rank is nonincreasing under
local maps in every power. Consequently $p$ maps the larger sublevel
onto the smaller one, and $j$ preserves asymptotic rank on the embedded
copy.

If the larger sublevel has components $C_1,\ldots,C_s$, the closures
of $p(C_1),\ldots,p(C_s)$ in the smaller sublevel are irreducible
closed sets covering it. Any irreducible component of the smaller
sublevel lies in one of those finitely many closed sets, and maximality
makes it equal to that set. The smaller sublevel therefore has at most
$s$ components. Its active-factor count is exactly $t_n$, so its
established lower bound gives the same bound in the original dimensions.
\end{proof}

When all $d_i\geq n$, this lower bound is $\binom{k}{3}+1$.
In particular, four factors of dimensions at least three have at least
five components for every $\sqrt[3]{62}\leq a<4$. The counts are lower
bounds, not complete lists of components. Their unboundedness varies
the tensor format; it makes no claim of infinitely many asymptotic-rank
values in a fixed format. The same section-and-retraction argument
also propagates reducibility when additional factors are appended
using fixed nonzero vectors.

\subsection{The grouped three-mode sublevel}
\label{sec:components-grouping}

\begin{proposition}\label{prop:components-grouping}
For every real $3\leq a<4$,
\begin{equation}\label{eq:components-grouped}
 Z_a(\C^9,\C^3,\C^2)
   =\{F:\rank F_{\C^9\mid\C^3\otimes\C^2}\leq3\}.
\end{equation}
This is the irreducible rank-at-most-three determinantal variety of
$9$-by-$6$ matrices.
\end{proposition}

\begin{proof}
Lemma~\ref{lem:flatten} gives the forward inclusion, using integrality
and $a<4$. Conversely, a tensor with this flattening rank at most three
has its first-mode support in a subspace of dimension at most three.
After compression and padding within that mode, it suffices to consider
format $(3,3,2)$.

Represent such a tensor by a pair of three-by-three matrices $(A,B)$.
The set where $A$ is invertible and $A^{-1}B$ has three distinct
eigenvalues is nonempty, as shown by $(I,\operatorname{diag}(0,1,2))$.
It is Zariski open: on $\det A\ne0$, distinct eigenvalues are the
nonvanishing of the characteristic discriminant. Thus it is dense in
the affine space of all pairs. Invertible left and right changes of
basis reduce every such pair to
\[
                  (I,\operatorname{diag}(\lambda_1,\lambda_2,\lambda_3)).
\]
The corresponding tensor is a sum of the three decomposable tensors
$e_i\otimes e_i\otimes(d_0+\lambda_i d_1)$. Its rank is at most three.
Taking Zariski closures proves that every tensor of format $(3,3,2)$
has border rank at most three and hence asymptotic rank at most three.
The latter implication also follows directly from closedness of $Z_3$:
that closed set contains every rank-at-most-three tensor and therefore
its closure. Local padding back to $(9,3,2)$ cannot increase the
asymptotic rank. Since $a\geq3$, this proves the reverse inclusion.

Every $9$-by-$6$ matrix of rank at most three is a product of a
$9$-by-$3$ matrix and a $3$-by-$6$ matrix. Hence the determinantal
variety is the image of an irreducible affine space under matrix
multiplication and is irreducible.
\end{proof}

Grouping any two of the three-dimensional modes of $(3,3,3,2)$ gives
$(9,3,2)$. Throughout $\sqrt[3]{62}\leq a<4$, the full four-mode
sublevel is reducible by Theorem~\ref{thm:components-count}, whereas
the full grouped three-mode sublevel is irreducible by
Proposition~\ref{prop:components-grouping}. This does not assert
surjectivity from the ungrouped sublevel onto the grouped one and
does not settle the general three-mode irreducibility question.

\subsection{Scope and the remaining three-mode questions}
\label{sec:components-scope}

The component lower bounds concern the full sublevels and allow the
format to vary. They neither classify all components nor produce
infinitely many asymptotic-rank values in a single fixed format.
The scalar-kernel argument uses three balanced flattenings and a
fourth factor containing simultaneous slices. The grouped example
shows explicitly why those constraints do not transfer to three
modes. A reducible linear section would also be insufficient to
prove reducibility of an ambient sublevel: the irreducible quadric
$xy=z^2$ has the reducible section $z=0$.

The full three-mode asymptotic-rank conjecture asserts
$\AR(T)=\max_i r_i(T)$, where $r_i(T)$ are the single-mode ranks,
after compression to those mode supports. Under this conjecture,
for every $a\geq0$ the sublevel is the subspace variety defined by
\[
 r_i(T)\leq\min\{\dim V_i,\lfloor a\rfloor\},\qquad i=1,2,3.
\]
It is irreducible: choose the three support spaces in the product
of the corresponding Grassmannians and then a tensor in their
tensor product. The resulting vector bundle has irreducible total
space and maps onto the subspace variety. When the bound is zero,
the variety is $\{0\}$. This conditional implication is already
recorded in \cite[Section~5]{CHNVZ2024}. Thus three-mode
reducibility would refute the full conjecture. It would not by
itself prove that the matrix-multiplication exponent $\omega$ is
greater than two, or refute the narrower conjecture for concise
tight cubic tensors \cite[Conjectures~3 and~8]{BjorklundKaski2023}.

\begin{proposition}\label{prop:components-finite-values}
In a fixed complex tensor format, finiteness of the set of
asymptotic-rank values is equivalent to a uniform bound on the
number of irreducible components of all its sublevels.
\end{proposition}

\begin{proof}
Finitely many values give finitely many sublevels, each a closed
affine variety with finitely many irreducible components.
Conversely, suppose all sublevels in an ambient affine space of
dimension $N$ have at most $C$ components. Associate to each one
its vector of component counts in dimensions $N,N-1,\ldots,0$.
For a strict inclusion of two closed subsets these vectors
strictly increase lexicographically. Indeed, until the first
dimension at which their components differ, all higher-dimensional
components coincide. An old component in the next dimension
cannot be absorbed by one of those common higher-dimensional
components, and containment in a new component of equal dimension
forces equality. Thus all old components of that dimension persist,
and at least one new component is added. If there were no differing
dimension, the two closed subsets would be equal.

There are only $\binom{C+N+1}{N+1}$ nonnegative integer vectors of
length $N+1$ with coordinate sum at most $C$. Consequently every
strict chain of sublevels has uniformly bounded finite length.
Distinct attained asymptotic-rank values give distinct, totally
ordered sublevels, so there are only finitely many such values.
\end{proof}

No uniform component bound in a fixed format is supplied by the
component-growth theorem. Determining the three-mode components in
the local product chart, finding a three-mode replacement for the
balanced-flattening argument, and controlling components over all
thresholds of a fixed format remain separate questions.

The arguments in this section are written proofs with the cited
closedness and asymptotic-speedup inputs. The finite certificate
described in Section~\ref{sec:verification} supports the original
permutation-tensor construction. It does not formalize the general
scalar cores, the component-growth theorem, the grouped
determinantal description, or the fixed-format proposition above.


\section{Exact verification and the finite formal layer}\label{sec:verification}
The finite calculations have two complementary forms: exact arithmetic
certificates for the explicit coefficients and matrices, and a Lean
proof of a proposition collecting their algebraic identities. All
calculations use integers, rational numbers, or Gaussian integers.
No floating-point tolerance, numerical rank decision, or recognition
of a decimal approximation is involved.

\subsection{Exact coefficient and matrix certificates}

The geometry calculation reconstructs all $1134$ rows of $L(P)$
from the appended determinants, using all $\binom94=126$ increasing
four-row choices for each relevant column and flattening. Rational
elimination gives rank $26$, and the identity $L(P)P=0$ identifies
its one-dimensional kernel. The same calculation checks the three
rank-three single-mode flattenings of $P$, the fourth-mode rank
two of $y$, and the three balanced ranks two of $y$.
A supplementary Gaussian change of basis takes $P$ to
$2\operatorname{cw}_2$; this identity is not needed in the geometric
argument of Section~\ref{sec:geometry}.

There are two independent coefficient calculations for the exchange.
One constructs the group, Walsh coefficients and sparse Laurent maps
directly, without reading the coefficient certificate. The other
contracts the specified $73$-to-$46$-dimensional maps and evaluates
all $46^3=97\,336$ ordered output coefficients. The source is written
using $88$ terms: $64$ Walsh summands and $24$ paired-$C_8$ terms.
These calculations check the cancellation of every pole and the
$270=216+54$ nonzero constant coefficients of
$P^{\boxtimes3}\oplus C_{18}$. They also check the ten-term
$\operatorname{cw}_8$ polynomial identity, the $19$ injective integer
weights for $C_{18}$, the zero weight sum on all $54$ support triples,
and all marginal counts. The counts are $18$ at the zero index
and $2$ at each other index, giving the rational probabilities
$1/3$ and $1/27$. The entropy evaluation and its asymptotic
consequence are the written calculation~\eqref{eq:entropy} and
Theorem~\ref{thm:tight}.

The three exact checkers reconstruct the expected tensors and reject
invalid input formats. Their runner binds the proof and certificate
files by hashes and compares registered output in ordinary execution,
optimized Python execution, and execution with a second hash seed.
One invocation includes nine positive executions and $48$ external
negative-control executions: each of $16$ specified corruptions must
be rejected in each of the three modes. The corruptions change, among
other quantities, the pole-cancellation sign, Fourier normalization,
a target coefficient, a degeneration scalar, the injectivity of tight
weights, or a witness rank. The coefficient checker also carries out
four corruption checks within each positive execution. Failure is
explicit and does not depend on assertions that optimization could
remove.

\subsection{The precise scope of the Lean certificate}

The formal proposition, with theorem name
\texttt{finiteCertificate}, combines the following statements.
It proves the complete enumeration of the minor rows and, for every
complex pair $T,S$, the equality between the appended determinant
and the corresponding coordinate of $L(T)S$. It proves $L(T)T=0$
for every complex $T$, the rank $26$ of $L(P)$, and
$L(P)S=0$ if and only if $S=cP$ for a complex scalar $c$.
The witness flattening ranks and the supplementary Gaussian
$P$-to-$2\operatorname{cw}_2$ change of basis are included.

For the exchange, the proposition proves the finite group and set
identities, Walsh orthogonality, agreement of the dense and sparse
contractions, and the permutation-coordinate identification of
$P^{\boxtimes3}$. It proves the Laurent coefficient identities at
every output triple, including vanishing in every degree below zero
or above three. The coefficient representation uses an integer
normalization, with constant coefficient $64$ times the specified
target; division by $64$ gives the normalization in
\eqref{eq:walsh}. The tight weights, marginal counts and integer
comparisons $27<62<64<81$ are included. The ten-term
$\operatorname{cw}_8$ identity is proved for every nonzero complex
parameter and all $9^3$ coordinate triples, with the required
constant value at zero of its regular polynomial expression.

The formal development uses Lean \texttt{v4.34.0-rc2} with a pinned
Mathlib revision. Its final theorem has precisely the logical
axioms \texttt{propext}, \texttt{Classical.choice} and
\texttt{Quot.sound}. Source and literal-data checks accompany a
clean rebuild and include controls for altered inputs, missing
proofs, altered build configuration or dependency manifests, and
unexpected axiom output.

This certificate does not formalize asymptotic rank or border rank,
Strassen's spectral or tight-support theorems, or the implication
from entropy to asymptotic subrank. It also does not formalize the
general balanced-rank-to-kernel equivalence of
Lemma~\ref{lem:appended}, the open-neighborhood argument, or the
proper relatively closed cover. The paired-$C_8$ to diagonal
$\operatorname{cw}_8$ basis change in Lemma~\ref{lem:cw-basis}
is a separate written bridge; the formalized Gaussian change for
$P$ does not replace it. None of these omitted statements is a
premise of the formal proposition. The asymptotic-rank bound and
the reducibility theorem follow from the finite identities together
with the stated classical inputs and the written arguments above.



\section*{Statements and Declarations}
\paragraph{Funding.}
The research of Wenqi Zhu was supported by the Engineering and Physical
Sciences Research Council under grant EP/Y028872/1, \emph{Mathematical
Foundations of Intelligence: An ``Erlangen Programme'' for AI}.

\paragraph{Competing interests.}
The authors have no competing interests to declare that are relevant to
the content of this article.

\paragraph{Data availability.}
The exact verification scripts, certificates and Lean~4 files accompanying
this paper are available at
\url{https://github.com/erenup/asymptotic-rank-sublevels}.

\paragraph{Use of generative AI and AI-assisted technologies.}
During the preparation of this work, the authors used AI tools (Codex and Claude Code) to generate candidate constructions and code, perform computations, draft mathematical arguments, and improve the readability of the text. After using these
tools, the authors reviewed and edited the content as needed and take full
responsibility for the content of the publication. The finite identities and the witness geometry are certified by exact arithmetic and by a Lean~4 finite certificate (Section~\ref{sec:verification}); the asymptotic-rank and Zariski arguments are given in full in the text.

\begingroup
\interlinepenalty=10000
\bibliographystyle{plainnat}
\bibliography{refs}
\endgroup
\end{document}
